\subsection{Hilbert Transform}

\subsubsection*{Motivation and Intuition}

The Hilbert transform is another linear transform used heavily in signal processing and communication theory. While the Fourier transform separates a signal into frequencies, the Hilbert transform is often used to create a phase-shifted version of a signal.

It is especially important in constructing analytic signals, where a real signal is paired with its Hilbert transform to form a complex signal.

\begin{conceptbox}
    The Hilbert transform of a function $x(t)$ is commonly denoted by $\hat{x}(t)$ or $H[x(t)]$ and is defined by the Cauchy principal value integral
    \[
        \hat{x}(t)
        =
        H[x(t)]
        =
        \frac{1}{\pi}\,\operatorname{P.V.}
        \int_{-\infty}^{\infty}
        \frac{x(\tau)}{t-\tau}\,d\tau.
    \]
\end{conceptbox}

The notation $\operatorname{P.V.}$ means principal value. This is necessary because the denominator becomes zero when $\tau=t$.

\subsubsection*{Physical Interpretation}

The Hilbert transform shifts the phase of frequency components. In the frequency domain, it corresponds to multiplication by $-i\,\operatorname{sgn}(\omega)$.

\begin{conceptbox}
    If
    \[
        X(\omega)=\mathcal{F}[x(t)],
    \]
    then the Fourier transform of the Hilbert transform is
    \[
        \mathcal{F}[\hat{x}(t)]
        =
        -i\,\operatorname{sgn}(\omega)X(\omega),
    \]
    where
    \[
        \operatorname{sgn}(\omega)
        =
        \begin{cases}
            1,  & \omega>0, \\
            0,  & \omega=0, \\
            -1, & \omega<0.
        \end{cases}
    \]
\end{conceptbox}

Thus, positive frequency components are multiplied by $-i$, corresponding to a phase shift of $-\pi/2$, while negative frequency components are multiplied by $i$, corresponding to a phase shift of $\pi/2$.

\subsubsection*{Key Properties}

\begin{conceptbox}
    \textbf{Linearity}
    \[
        H[c_1x(t)+c_2y(t)]
        =
        c_1H[x(t)]+c_2H[y(t)].
    \]
\end{conceptbox}

\begin{conceptbox}
    \textbf{Frequency-Domain Form}
    \[
        \widehat{X}(\omega)
        =
        -i\,\operatorname{sgn}(\omega)X(\omega).
    \]
\end{conceptbox}

\begin{conceptbox}
    \textbf{Analytic Signal}
    The analytic signal associated with $x(t)$ is
    \[
        z(t)=x(t)+i\hat{x}(t).
    \]
    The real part is the original signal, and the imaginary part is its Hilbert transform.
\end{conceptbox}

\subsubsection*{Hilbert Transforms of Common Functions}

\begin{conceptbox}
    Using the convention
    \[
        \mathcal{F}[H[x(t)]]
        =
        -i\,\operatorname{sgn}(\omega)X(\omega),
    \]
    the Hilbert transform shifts positive-frequency components by $-\pi/2$ and negative-frequency components by $\pi/2$.

    \[
        \begin{array}{c|c}
            \textbf{Function } x(t) & \textbf{Hilbert Transform } H[x(t)]            \\ \hline \\[1pt]
            1                       & 0                                              \\[6pt]
            \cos(\omega_0t)         & \sin(\omega_0t)                                \\[6pt]
            \sin(\omega_0t)         & -\cos(\omega_0t)                               \\[6pt]
            e^{i\omega_0t}          & -i\,\operatorname{sgn}(\omega_0)e^{i\omega_0t} \\[8pt]
            \delta(t)               & \dfrac{1}{\pi t}                               \\[10pt]
            \dfrac{1}{\pi t}        & -\delta(t)
        \end{array}
    \]
\end{conceptbox}

\begin{conceptbox}
    For a real signal $x(t)$, the analytic signal is
    \[
        z(t)=x(t)+iH[x(t)].
    \]

    Common analytic signal examples are:
    \[
        \begin{array}{c|c}
            \textbf{Real Signal } x(t) & \textbf{Analytic Signal } z(t) \\ \hline
            \cos(\omega_0t)            & e^{i\omega_0t}                 \\[6pt]
            \sin(\omega_0t)            & -ie^{i\omega_0t}
        \end{array}
    \]
\end{conceptbox}

\begin{noteBox}
    The Hilbert transform is useful when the phase structure of a signal matters. In communication systems, it helps form analytic signals, amplitude envelopes, instantaneous phase, and single-sideband modulated signals.
\end{noteBox}

\subsubsection*{Worked Examples}

\begin{examplebox}
    \textbf{Example 1: Hilbert transform of a cosine signal}

    Let
    \[
        x(t)=\cos(\omega_0t).
    \]

    Using Euler's formula,
    \[
        \cos(\omega_0t)
        =
        \frac{1}{2}
        \left(e^{i\omega_0t}+e^{-i\omega_0t}\right).
    \]

    The Fourier transform contains impulses at $\omega=\omega_0$ and $\omega=-\omega_0$:
    \[
        X(\omega)=\pi\left[\delta(\omega-\omega_0)+\delta(\omega+\omega_0)\right].
    \]

    Apply the Hilbert transform frequency rule:
    \[
        \widehat{X}(\omega)
        =
        -i\,\operatorname{sgn}(\omega)X(\omega).
    \]

    At $\omega=\omega_0>0$,
    \[
        -i\,\operatorname{sgn}(\omega_0)=-i.
    \]

    At $\omega=-\omega_0<0$,
    \[
        -i\,\operatorname{sgn}(-\omega_0)=i.
    \]

    Thus,
    \[
        \widehat{X}(\omega)
        =
        -i\pi\delta(\omega-\omega_0)
        +
        i\pi\delta(\omega+\omega_0).
    \]

    This is the Fourier transform of
    \[
        \sin(\omega_0t).
    \]

    Therefore,
    \[
        \boxed{
            H[\cos(\omega_0t)] = \sin(\omega_0t)
        }.
    \]
\end{examplebox}

\begin{examplebox}
    \textbf{Example 2: Hilbert transform of a sine signal}

    Let
    \[
        x(t)=\sin(\omega_0t).
    \]

    Using Euler's formula,
    \[
        \sin(\omega_0t)
        =
        \frac{1}{2i}
        \left(e^{i\omega_0t}-e^{-i\omega_0t}\right).
    \]

    Its Hilbert transform shifts the positive-frequency component by $-\pi/2$ and the negative-frequency component by $\pi/2$.

    Using the known phase-shift relation,
    \[
        H[\sin(\omega_0t)] = -\cos(\omega_0t).
    \]

    Therefore,
    \[
        \boxed{
            H[\sin(\omega_0t)] = -\cos(\omega_0t)
        }.
    \]
\end{examplebox}

\begin{examplebox}
    \textbf{Example 3: Analytic signal for $x(t)=\cos(\omega_0t)$}

    From Example 1,
    \[
        H[\cos(\omega_0t)] = \sin(\omega_0t).
    \]

    The analytic signal is
    \[
        z(t)=x(t)+i\hat{x}(t).
    \]

    Substitute:
    \[
        z(t)=\cos(\omega_0t)+i\sin(\omega_0t).
    \]

    Using Euler's formula,
    \[
        z(t)=e^{i\omega_0t}.
    \]

    Thus, the analytic signal associated with $\cos(\omega_0t)$ is
    \[
        \boxed{
            z(t)=e^{i\omega_0t}
        }.
    \]
\end{examplebox}

\subsubsection*{Engineering Insight}

\begin{noteBox}
    The Hilbert transform is used in communication systems to form analytic signals, determine instantaneous phase, compute instantaneous frequency, and analyze amplitude envelopes. It is also used in single-sideband modulation, radar signal processing, and spectral-response analysis.
\end{noteBox}

\subsubsection*{Common Mistakes}

\begin{conceptbox}
    Common mistakes in Hilbert transforms include:
    \[
        \begin{aligned}
             & \text{1. Ignoring the principal value nature of the integral.}                             \\
             & \text{2. Treating the Hilbert transform as an ordinary time delay.}                        \\
             & \text{3. Forgetting that it changes phase but not amplitude.}                              \\
             & \text{4. Confusing the analytic signal } z(t) \text{ with the original real signal } x(t). \\
             & \text{5. Applying the transform without considering whether the signal is suitable.}
        \end{aligned}
    \]
\end{conceptbox}